\section{Related Work}

% Framing paragraph: narrative disclosures and market reactions
This section relates the paper to the literature on narrative disclosures and capital market reactions. Prior research establishes that textual information is relevant for asset prices. \citet{tetlock2007giving} shows that pessimistic language in financial news predicts market outcomes, while \citet{tetlock2008morethanwords} show that linguistic features of firm-specific news are informative about firm fundamentals and stock returns. In issuer-authored disclosure settings, narrative text also conveys economically meaningful information beyond numerical accounting data. In the context of regulatory filings, large-sample evidence shows that investors respond to narrative content in MD\&A at the time of disclosure. For example, changes in MD\&A text from year to year are positively associated with filing-date abnormal returns, suggesting that narrative updates convey incremental information to equity markets \citep{brown2011mda}. Similarly, studies of MD\&A tone changes document abnormal return responses around filing events even after controlling for earnings surprises and accruals, highlighting the independent role of managerial narratives in price formation \citep{feldman2010tone}. Together, these findings establish MD\&A disclosures as a relevant setting for studying how markets incorporate narrative information. Relatedly, \citet{henry2016measuring} show that in disclosure settings, inferences about tone informativeness depend importantly on how tone is measured.

% Measurement of textual sentiment in finance
A substantial body of work studies how sentiment and related textual attributes should be measured in financial documents. Early evidence shows that domain-specific dictionaries materially outperform generic lexicons in financial disclosure settings because general-language wordlists misclassify many finance-specific usages \citep{loughran2011liability}. Closely related, \citet{henry2016measuring} directly compare alternative methodologies for measuring disclosure tone and show that domain-specific wordlists outperform generic ones in explaining market reactions around disclosure dates. They further find that simple finance-specific word-frequency measures are generally as informative as more complex machine-learning-based alternatives, while offering greater transparency, ease of implementation, and replicability. Consistent with this broader measurement perspective, methodological surveys emphasize that dictionary-based proxies remain attractive for empirical research but are sensitive to construction choices and may contain nontrivial measurement error \citep{loughranmcdonald2016survey}. Taken together, this literature suggests that inference about disclosure informativeness depends importantly on how sentiment is quantified.

% Broad use of text as data in finance
More broadly, recent text-as-data research in economics shows that corporate narratives can be used to measure several economically meaningful constructs, including exposure, risk, and sentiment. It further suggests that the choice of text-analysis method should be guided by the research objective, the structure of the corpus, and the tradeoff between flexibility and interpretability \citep{hoberg2016text, hassan2025text}. For example, \citet{hoberg2016text} use text-based analysis of firms’ 10-K product descriptions to construct time-varying measures of product market similarity and show that these text-derived measures outperform traditional industry classifications in explaining competition, profitability, and firm behavior. These studies underscore that narrative disclosures can be informative about firm fundamentals and strategic positioning, motivating subsequent work that seeks to extract specific informational dimensions—such as sentiment—from financial texts. Building on this broader ``text-as-data'' perspective, we focus on a specific informational dimension—tone in Japanese EDINET MD\&A—and examine whether alternative sentiment constructions (dictionary-based versus LLM-based) lead to different inferences about market reactions around filing dates.

% Cross-linguistic evidence and international settings
While finance-specific dictionaries provide a dominant framework for sentiment analysis in English-language disclosures, extending these methods to non-English settings introduces additional linguistic and institutional challenges. Cross-country studies document that naive translation of English financial sentiment dictionaries leads to substantial measurement error due to language-specific features such as morphology and compounding, often requiring substantial adaptation or reconstruction of lexicons from native-language corpora \citep{bannier2018german, du2022chinese, kim2021korean}. This evidence suggests that dictionary-based sentiment measures are strongly shaped by both language and domain, raising further concerns about measurement reliability in international disclosure regimes. This concern is conceptually consistent with evidence from English-language disclosure settings that financial reporting constitutes a specialized sublanguage, so sentiment measures that ignore domain-specific usage are prone to misclassification even before cross-lingual translation issues are introduced \citep{henry2016measuring}.

\subsection{Transformer-based encoders: BERT and domain adaptation}

Transformer architectures enabled substantial gains in text classification by pretraining on large corpora and fine-tuning on downstream tasks. \citet{devlin2019bert} introduce BERT, a bidirectional Transformer encoder trained with masked language modeling, which can be fine-tuned efficiently for classification problems such as sentiment prediction. In financial applications, domain-adapted variants (FinBERT) further improve performance by continuing pretraining on financial corpora before task-specific fine-tuning \citep{yang2020finbert}. These models provide a natural baseline for supervised sentiment extraction in financial disclosures, including non-English settings when suitable pretrained models or domain corpora are available.

\subsection{Autoregressive Large Language Models}

Beyond encoder-based Transformer models such as BERT, autoregressive Transformer models pretrained to predict the next token (the GPT family) have demonstrated strong performance across many NLP tasks and can be applied to document-level inference via prompting \citep{radford2018improving}. Modern chat-oriented systems further apply instruction tuning and reinforcement learning from human feedback to improve alignment with user intent \citep{ouyang2022instructgpt}, while later generations such as GPT-4 demonstrate substantially improved benchmark performance \citep{openai2023gpt4}. In financial applications, these models can be used directly for document-level sentiment inference via prompting, but their post-training alignment may also shape outputs (e.g., systematic positivity), motivating careful validation and robustness checks when applied to economic inference.

% Japanese LLM evidence: positioned as predictability, not event studies
Recent accounting evidence also shows that contextualized LLM-based measures can substantially outperform dictionary and non-LLM text measures in explaining market reactions to earnings announcements, suggesting that word context may materially reduce measurement error in disclosure-based inference \citep{siano2025news}. Consistent with this broader shift toward contextual methods, \citet{okada2025words} apply large language models directly to Japanese corporate disclosures and document return predictability from LLM-derived sentiment measures, whereas dictionary-based metrics exhibit limited explanatory power. However, this evidence primarily evaluates portfolio-level predictability rather than event-window reactions around disclosure dates. It therefore stands in some tension with evidence from accounting disclosure settings that domain-specific and comparatively transparent tone measures can outperform broader or more complex alternatives \citep{henry2016measuring}. This leaves open whether contextual LLM-based sentiment measures continue to dominate in structured regulatory filings, particularly in non-English settings where both domain specificity and language adaptation are central measurement concerns.


% Gap statement: event-study and inference robustness
Despite extensive evidence that narrative disclosures affect asset prices, important uncertainty remains regarding which sentiment extraction methods are most informative in disclosure-based event studies. Prior accounting evidence shows that in specialized disclosure settings, domain-specific and transparent tone measures can outperform broader or more complex alternatives \citep{henry2016measuring}, whereas recent evidence from English earnings-announcement disclosures finds that contextualized LLM-based measures can substantially outperform dictionary-based and non-LLM methods in explaining short-window market reactions \citep{siano2025news}. These literatures have not been fully reconciled in the setting of structured Japanese regulatory filings. This issue is particularly salient in non-English disclosure regimes, where linguistic structure and institutional formatting may exacerbate measurement error. Accordingly, it remains unclear whether alternative sentiment extraction approaches yield systematically different abnormal return responses around disclosure events, and thus whether method choice materially affects economic inference in standard event-study settings.
